\section{Introduction}\label{sec:intro}

In the ReAct pattern~\cite{yao2023react}, an LLM agent interleaves reasoning with tool calls such as searching
documents, querying a database or sending e-mail. Whatever a tool returns enters the model's context directly. In an
indirect prompt injection (IPI) attack~\cite{greshake2023indirect}, the attacker plants instructions in third-party
content that the agent will read. If the model follows them, the agent may send data to an unknown destination, leak a
confidential value or alter data that the user asked it to send.

Agent benchmarks such as BIPIA~\cite{yi2025bipia}, InjecAgent~\cite{zhan2024injecagent} and
AgentDojo~\cite{debenedetti2024agentdojo} are almost entirely in English. Vietnamese text is often written without
diacritics and mixed with English terms, while injection classifiers are mainly trained on English. Prompt-level
defences such as spotlighting~\cite{hines2024spotlighting} depend on whether the model complies. This paper therefore
places the enforcement rules on the host side, between the model and the tools, where they do not depend on model
compliance. The reminder in the prompt is only a supporting layer.

The paper makes four contributions.
\begin{enumerate}
\item \textbf{A mechanism-based threat model.} Two independent labels (trust and sensitivity), four harm predicates,
  six attack mechanisms and two benign counterparts for measuring false blocks (Section~\ref{sec:threat}).
\item \textbf{The \ours{} defence.} Rules on destinations and on the flow of confidential values are enforced in the
  Tool Broker and combined with a prompt reminder and denial feedback, without fine-tuning. Rules and reminder
  complement each other across attack mechanisms. A proposition states what the rules guarantee and under which
  assumptions, and every remaining bypass either violates one of these assumptions or is a parameter-tampering attack
  (Sections~\ref{sec:defence} and~\ref{sec:discussion}).
\item \textbf{VN-AgentInject and AgentDojo-VI.} VN-AgentInject is evaluated under a sealed protocol. The test sets,
  defence configurations and hypotheses are locked by hash before the run, and each test set is run only once.
  AgentDojo-VI is a Vietnamese translation of AgentDojo used as an external check (Section~\ref{sec:protocol}).
  The data, code and seal manifests are publicly released.
\item \textbf{Empirical results.} Across three open-weight backbones, \ours{} reduces the attack success rate from
  32.9\% to 1.2\% on \tOne{} and from 32.4\% to 3.8\% on \tTwo, a set written after \ours{} was frozen. Benign
  utility drops by about 5 percentage points. The results are compared with six baselines and Gemini 3 Flash on Dev,
  and with AgentDojo in English and Vietnamese (Section~\ref{sec:results}).
\end{enumerate}
